% Regression: xyz_test — imported current-behavior fixture from the handoff corpus.
% Formula: (a) ONE tenkzlattice with sheets=2 + pairing reproduces the
% xyz_test — acceptance for the projected-basis-vector lattice frame.
%   (a) ONE tenkzlattice with sheets=2 + pairing reproduces the
%       tenkzplanes double layer (strict byte-compare of the two lone
%       pictures lives in xyz_a_lattice.tex / xyz_a_planes.tex);
%   (b) per-sheet decoration: a bra-sheet \tnedge (r,c,0)-(r',c',0), a
%       ket-only region, an inter-sheet tie — inside tenkzplanes, whose
%       unaddressed cells still mean the ket sheet;
%   (c) an expert custom projection via col vector / row vector;
%   (d) the existing tenkzplanes spellings, unchanged.
% Expected warnings: plane-rise-low at (a)/(b') only — the lattice
% spelling states rise=0.45 through the guarded public key, while the
% tenkzplanes preset installs the same planes-local ratio ungated; the
% guard prices single-sheet physical legs, which a double layer has none
% of, so it warns and draws (documented key behaviour, kept).
\documentclass[varwidth=20cm, border=6pt]{standalone}
\usepackage{amsmath, amssymb}
\usepackage{tenkz}
\begin{document}

% ---- (a) one lattice, two sheets = the double layer --------------------
% sheets=2 stacks two oblique sheets along S = (0, sep) with the derived
% sep = (rows-1) x rise + gap; pairing draws the physical leg along S at
% every column.  With the planes-local ratios stated explicitly, this IS
% the tenkzplanes picture below it (bra = sheet 0, ket = sheet 1 above).
\textbf{(a) tenkzlattice, sheets=2 + pairing}\par\smallskip
\begin{tenkzlattice}[rows=3, cols=3, sheets=2, pairing, frame=oblique,
                     plane slant=0.40, plane rise=0.45]
\end{tenkzlattice}

\medskip
\textbf{(a) the same picture as tenkzplanes}\par\smallskip
\begin{tenkzplanes}[rows=3, cols=3]
\end{tenkzplanes}

\medskip
% ---- (b) per-sheet decoration inside tenkzplanes -----------------------
% Unaddressed cells keep the Phase-1 meaning (ket sheet): the region
% $K$ lives on the ket sheet ONLY.  Sheet 0 addresses the bra sheet: the
% distinguished bra string (1,1,0)-(1,3,0) is the t2_condensanyon
% blocker, and the tie (1,2,1)-(1,2,0) marks one column's contraction.
\textbf{(b) tenkzplanes: bra-sheet edge, ket-only region, tie}%
\par\smallskip
\begin{tenkzplanes}[rows=3, cols=3]
  \tnregion[slot=selected, label=$K$]{(2-3,2-3)}
  \tnedge{(1,1,0)-(1,3,0)}
  \tnedge{(1,2,1)-(1,2,0)}
\end{tenkzplanes}

\medskip
% (b') the same addressing through the tenkzlattice spelling, plus a
% bra-only region: sheet coordinates work identically in both
% environments (here the default sheet is 0, so the ket region says ,1)
\textbf{(b') tenkzlattice, sheets=2: per-sheet regions}\par\smallskip
\begin{tenkzlattice}[rows=3, cols=3, sheets=2, pairing, frame=oblique,
                     plane slant=0.40, plane rise=0.45]
  \tnregion[slot=selected, label=$K$]{(2-3,2-3,1)}
  \tnregion[slot=secondary, label=$B$, label at=south west]{(1-2,1-2,0)}
  \tnedge{(1,1,0)-(1,3,0)}
\end{tenkzlattice}

\medskip
% ---- (c) expert projection ---------------------------------------------
% col vector / row vector state the projected basis directly: a taller,
% steeper frame than any preset — C = (9mm, 1.5mm) tilts the rows up,
% R = (-3mm, -10mm) drops each row a full 10mm at a 3mm westward lean.
% The frame counts as sheared (R_x != 0), so the region label flips to
% the outside corner; bonds, tracer and boundary legs follow the basis
% while dots stay round and the label upright (nodes never shear).
\textbf{(c) expert frame: col vector=\{9mm,1.5mm\},
  row vector=\{-3mm,-10mm\}}\par\smallskip
\begin{tenkzlattice}[rows=3, cols=3, col vector={9mm,1.5mm},
                     row vector={-3mm,-10mm}, boundary legs]
  \tnregion[slot=selected, label=$R$]{(1-2,2-3)}
\end{tenkzlattice}

\medskip
% ---- (d) existing tenkzplanes spellings, unchanged ---------------------
\textbf{(d) tenkzplanes, derived sep + open columns (Phase-1 spelling)}%
\par\smallskip
\begin{tenkzplanes}[rows=3, cols=3, sheets={ket, bra}]
\end{tenkzplanes}

\medskip
\begin{tenkzplanes}[rows=3, cols=3, open={(1,2),(3,3)}]
\end{tenkzplanes}

\end{document}
